\documentclass[12pt]{article}

\usepackage{graphicx}			% Use this package to include images
\usepackage{amsmath}			% A library of many standard math expressions
\usepackage[margin=1in]{geometry}% Sets 1in margins. 
\usepackage{fancyhdr}			% Creates headers and footers
\usepackage{enumerate}          %These two package give custom labels to a list
\usepackage[shortlabels]{enumitem}
\usepackage{matlab-prettifier}
\usepackage{tikz}
\usepackage[utf8]{inputenc}
\usepackage{pgfplots} 
\usepackage{pgfgantt}
\usepackage{pdflscape}
\usepackage{array,multirow}
\usepackage{tabularx,booktabs}
\newcolumntype{Y}{>{\centering\arraybackslash}X}
\usetikzlibrary{shapes.geometric}
\usetikzlibrary{arrows.meta,arrows}
\pgfplotsset{compat=newest} 
\pgfplotsset{plot coordinates/math parser=false}
\newlength\fwidth
\newlength\fheight
\usepackage{float,graphicx}
\pagestyle{fancy}
\fancyhead[l]{}
\fancyhead[c]{Experiment 1}
\fancyhead[r]{}
\fancyfoot[c]{\thepage}
\setlength{\headheight}{20pt}
\begin{document}
\author{}
\date{}
\title{Stefan's Law}
\maketitle
\section*{Aim}
To verify that a body at temperature T radiates heat at a rate proportional to the fourth power of T and calculate the constant of proportionality.

\section*{Theory}
For a body at temperature $T$ and ambient temperature $T_a$, the rate of heat exchange is given by the expression
\[\dot{Q}=\epsilon \sigma A(T^4 -T_a^4)\]
where $\sigma$ is Stefan's constant, $\epsilon$ is the emissivity and $A$ is the surface area of the body.
In the experimental setup, a tungsten filament is connected to a battery and emits radiation.
When the system archives thermal equilibrium, the electrical power consumed is equal to the rate of heat radiated.
It is also given that the resistance of the filament depends on the temperature according to the given equation
\[R = R_0(1+\alpha T +\beta T^2)\]
\[\beta T^2 + \alpha T + (1-\frac{R}{R_0})=0\]

\begin{align*}
    T &= \frac{-\alpha + \sqrt{\alpha ^2 -4\beta (1-\frac{R}{R_0})} }{2\beta} \\
    T &= \frac{-5.21 + \sqrt{27.1441-2.88(1-\frac{R}{1.31})}}{14.4}\times 10^4
\end{align*}
\section*{Analysis}
The actual dependence of temperature is given by
\[\frac{dR}{R}=c dT\]
where c is a constant. Solving this equation gives
\[R=R_0e^{c(T-T_0)}\]

The given relation of the dependence of temperature is the taylor series approximation
This is a good approximation but has a few problems that are immediately apparent.
First we see the fisrt few calculated values of resistance by the quadratic relation.
\begin{table}[!hbt]
    \centering
    \begin{tabular}{|c|c|c|c|}
    \hline
        Voltage($V$) & Current($I$) & Resistance($R$) & Temperature($T$) \\ \hline
        0.03 & 0.01 & 3.00 & 239.61 \\ \hline
        0.07 & 0.02 & 3.50 & 307.70 \\ \hline
        0.11 & 0.03 & 3.67 & 330.14 \\ \hline
        0.16 & 0.04 & 4.00 & 374.64 \\ \hline
        0.22 & 0.05 & 4.40 & 427.39 \\ \hline
        0.34 & 0.06 & 5.67 & 590.07 \\ \hline
        0.48 & 0.07 & 6.86 & 737.43 \\ \hline
        0.59 & 0.08 & 7.38 & 800.00 \\ \hline
    \end{tabular}
    \caption{Values of $R$ and $T$ for lower values}
\end{table}

Notice that the first value of the calculated resistance is more than 50$K$ below the ambient temperature ($296K$) which cannot be possible.
To see what is going wrong, we look at the graph below.
\begin{figure}[!ht]
    \centering
    \setlength\fwidth{0.62\textwidth}
    \setlength\fheight{0.5\textwidth}
    % This file was created by matlab2tikz.
%
\definecolor{mycolor1}{rgb}{0.00000,0.44700,0.74100}%
\definecolor{mycolor2}{rgb}{0.85000,0.32500,0.09800}%
\definecolor{mycolor3}{rgb}{0.12941,0.12941,0.12941}%
%
\begin{tikzpicture}

\begin{axis}[%
width=0.951\fwidth,
height=\fheight,
at={(0\fwidth,0\fheight)},
scale only axis,
xmin=0,
xmax=2500,
xlabel style={font=\color{mycolor3}},
xlabel={Temperature},
ymin=0.1,
ymax=25,
ylabel style={font=\color{mycolor3}},
ylabel={Resistance},
axis background/.style={fill=white},
xmajorgrids,
ymajorgrids,
legend style={at={(0.03,0.97)}, anchor=north west, legend cell align=left, align=left}
]
\addplot [color=mycolor1]
  table[row sep=crcr]{%
-500	-2.48495175\\
-490	-2.4189678607\\
-480	-2.3528323328\\
-470	-2.2865451663\\
-460	-2.2201063612\\
-450	-2.1535159175\\
-440	-2.0867738352\\
-430	-2.0198801143\\
-420	-1.9528347548\\
-410	-1.8856377567\\
-400	-1.81828912\\
-390	-1.7507888447\\
-380	-1.6831369308\\
-370	-1.6153333783\\
-360	-1.5473781872\\
-350	-1.4792713575\\
-340	-1.4110128892\\
-330	-1.3426027823\\
-320	-1.2740410368\\
-310	-1.2053276527\\
-300	-1.13646263\\
-290	-1.0674459687\\
-280	-0.9982776688\\
-270	-0.9289577303\\
-260	-0.8594861532\\
-250	-0.7898629375\\
-240	-0.7200880832\\
-230	-0.6501615903\\
-220	-0.5800834588\\
-210	-0.5098536887\\
-200	-0.43947228\\
-190	-0.3689392327\\
-180	-0.2982545468\\
-170	-0.2274182223\\
-160	-0.1564302592\\
-150	-0.0852906574999999\\
-140	-0.0139994172\\
-130	0.0574434617\\
-120	0.1290379792\\
-110	0.2007841353\\
-100	0.27268193\\
-90	0.3447313633\\
-80	0.4169324352\\
-70	0.4892851457\\
-60	0.5617894948\\
-50	0.6344454825\\
-40	0.7072531088\\
-30	0.7802123737\\
-20	0.8533232772\\
-10	0.9265858193\\
0	1\\
10	1.0735658193\\
20	1.1472832772\\
30	1.2211523737\\
40	1.2951731088\\
50	1.3693454825\\
60	1.4436694948\\
70	1.5181451457\\
80	1.5927724352\\
90	1.6675513633\\
100	1.74248193\\
110	1.8175641353\\
120	1.8927979792\\
130	1.9681834617\\
140	2.0437205828\\
150	2.1194093425\\
160	2.1952497408\\
170	2.2712417777\\
180	2.3473854532\\
190	2.4236807673\\
200	2.50012772\\
210	2.5767263113\\
220	2.6534765412\\
230	2.7303784097\\
240	2.8074319168\\
250	2.8846370625\\
260	2.9619938468\\
270	3.0395022697\\
280	3.1171623312\\
290	3.1949740313\\
300	3.27293737\\
310	3.3510523473\\
320	3.4293189632\\
330	3.5077372177\\
340	3.5863071108\\
350	3.6650286425\\
360	3.7439018128\\
370	3.8229266217\\
380	3.9021030692\\
390	3.9814311553\\
400	4.06091088\\
410	4.1405422433\\
420	4.2203252452\\
430	4.3002598857\\
440	4.3803461648\\
450	4.4605840825\\
460	4.5409736388\\
470	4.6215148337\\
480	4.7022076672\\
490	4.7830521393\\
500	4.86404825\\
510	4.9451959993\\
520	5.0264953872\\
530	5.1079464137\\
540	5.1895490788\\
550	5.2713033825\\
560	5.3532093248\\
570	5.4352669057\\
580	5.5174761252\\
590	5.5998369833\\
600	5.68234948\\
610	5.7650136153\\
620	5.8478293892\\
630	5.9307968017\\
640	6.0139158528\\
650	6.0971865425\\
660	6.1806088708\\
670	6.2641828377\\
680	6.3479084432\\
690	6.4317856873\\
700	6.51581457\\
710	6.5999950913\\
720	6.6843272512\\
730	6.7688110497\\
740	6.8534464868\\
750	6.9382335625\\
760	7.0231722768\\
770	7.1082626297\\
780	7.1935046212\\
790	7.2788982513\\
800	7.36444352\\
810	7.4501404273\\
820	7.5359889732\\
830	7.6219891577\\
840	7.7081409808\\
850	7.7944444425\\
860	7.8808995428\\
870	7.9675062817\\
880	8.0542646592\\
890	8.1411746753\\
900	8.22823633\\
910	8.3154496233\\
920	8.4028145552\\
930	8.4903311257\\
940	8.5779993348\\
950	8.6658191825\\
960	8.7537906688\\
970	8.8419137937\\
980	8.9301885572\\
990	9.0186149593\\
1000	9.107193\\
1010	9.1959226793\\
1020	9.2848039972\\
1030	9.3738369537\\
1040	9.4630215488\\
1050	9.5523577825\\
1060	9.6418456548\\
1070	9.7314851657\\
1080	9.8212763152\\
1090	9.9112191033\\
1100	10.00131353\\
1110	10.0915595953\\
1120	10.1819572992\\
1130	10.2725066417\\
1140	10.3632076228\\
1150	10.4540602425\\
1160	10.5450645008\\
1170	10.6362203977\\
1180	10.7275279332\\
1190	10.8189871073\\
1200	10.91059792\\
1210	11.0023603713\\
1220	11.0942744612\\
1230	11.1863401897\\
1240	11.2785575568\\
1250	11.3709265625\\
1260	11.4634472068\\
1270	11.5561194897\\
1280	11.6489434112\\
1290	11.7419189713\\
1300	11.83504617\\
1310	11.9283250073\\
1320	12.0217554832\\
1330	12.1153375977\\
1340	12.2090713508\\
1350	12.3029567425\\
1360	12.3969937728\\
1370	12.4911824417\\
1380	12.5855227492\\
1390	12.6800146953\\
1400	12.77465828\\
1410	12.8694535033\\
1420	12.9644003652\\
1430	13.0594988657\\
1440	13.1547490048\\
1450	13.2501507825\\
1460	13.3457041988\\
1470	13.4414092537\\
1480	13.5372659472\\
1490	13.6332742793\\
1500	13.72943425\\
1510	13.8257458593\\
1520	13.9222091072\\
1530	14.0188239937\\
1540	14.1155905188\\
1550	14.2125086825\\
1560	14.3095784848\\
1570	14.4067999257\\
1580	14.5041730052\\
1590	14.6016977233\\
1600	14.69937408\\
1610	14.7972020753\\
1620	14.8951817092\\
1630	14.9933129817\\
1640	15.0915958928\\
1650	15.1900304425\\
1660	15.2886166308\\
1670	15.3873544577\\
1680	15.4862439232\\
1690	15.5852850273\\
1700	15.68447777\\
1710	15.7838221513\\
1720	15.8833181712\\
1730	15.9829658297\\
1740	16.0827651268\\
1750	16.1827160625\\
1760	16.2828186368\\
1770	16.3830728497\\
1780	16.4834787012\\
1790	16.5840361913\\
1800	16.68474532\\
1810	16.7856060873\\
1820	16.8866184932\\
1830	16.9877825377\\
1840	17.0890982208\\
1850	17.1905655425\\
1860	17.2921845028\\
1870	17.3939551017\\
1880	17.4958773392\\
1890	17.5979512153\\
1900	17.70017673\\
1910	17.8025538833\\
1920	17.9050826752\\
1930	18.0077631057\\
1940	18.1105951748\\
1950	18.2135788825\\
1960	18.3167142288\\
1970	18.4200012137\\
1980	18.5234398372\\
1990	18.6270300993\\
2000	18.730772\\
2010	18.8346655393\\
2020	18.9387107172\\
2030	19.0429075337\\
2040	19.1472559888\\
2050	19.2517560825\\
2060	19.3564078148\\
2070	19.4612111857\\
2080	19.5661661952\\
2090	19.6712728433\\
2100	19.77653113\\
2110	19.8819410553\\
2120	19.9875026192\\
2130	20.0932158217\\
2140	20.1990806628\\
2150	20.3050971425\\
2160	20.4112652608\\
2170	20.5175850177\\
2180	20.6240564132\\
2190	20.7306794473\\
2200	20.83745412\\
2210	20.9443804313\\
2220	21.0514583812\\
2230	21.1586879697\\
2240	21.2660691968\\
2250	21.3736020625\\
2260	21.4812865668\\
2270	21.5891227097\\
2280	21.6971104912\\
2290	21.8052499113\\
2300	21.91354097\\
2310	22.0219836673\\
2320	22.1305780032\\
2330	22.2393239777\\
2340	22.3482215908\\
2350	22.4572708425\\
2360	22.5664717328\\
2370	22.6758242617\\
2380	22.7853284292\\
2390	22.8949842353\\
2400	23.00479168\\
2410	23.1147507633\\
2420	23.2248614852\\
2430	23.3351238457\\
2440	23.4455378448\\
2450	23.5561034825\\
2460	23.6668207588\\
2470	23.7776896737\\
2480	23.8887102272\\
2490	23.9998824193\\
2500	24.11120625\\
2510	24.2226817193\\
2520	24.3343088272\\
2530	24.4460875737\\
2540	24.5580179588\\
2550	24.6700999825\\
2560	24.7823336448\\
2570	24.8947189457\\
2580	25.0072558852\\
2590	25.1199444633\\
2600	25.23278468\\
2610	25.3457765353\\
2620	25.4589200292\\
2630	25.5722151617\\
2640	25.6856619328\\
2650	25.7992603425\\
2660	25.9130103908\\
2670	26.0269120777\\
2680	26.1409654032\\
2690	26.2551703673\\
2700	26.36952697\\
2710	26.4840352113\\
2720	26.5986950912\\
2730	26.7135066097\\
2740	26.8284697668\\
2750	26.9435845625\\
2760	27.0588509968\\
2770	27.1742690697\\
2780	27.2898387812\\
2790	27.4055601313\\
2800	27.52143312\\
2810	27.6374577473\\
2820	27.7536340132\\
2830	27.8699619177\\
2840	27.9864414608\\
2850	28.1030726425\\
2860	28.2198554628\\
2870	28.3367899217\\
2880	28.4538760192\\
2890	28.5711137553\\
2900	28.68850313\\
2910	28.8060441433\\
2920	28.9237367952\\
2930	29.0415810857\\
2940	29.1595770148\\
2950	29.2777245825\\
2960	29.3960237888\\
2970	29.5144746337\\
2980	29.6330771172\\
2990	29.7518312393\\
3000	29.870737\\
};
\addlegendentry{Quadratic}

\addplot [color=mycolor2, dashed]
  table[row sep=crcr]{%
-500	-8.86609156836105\\
-490	-8.75539840395572\\
-480	-8.64470187562285\\
-470	-8.53400198325956\\
-460	-8.42329872676419\\
-450	-8.31259210603468\\
-440	-8.20188212096775\\
-430	-8.09116877146216\\
-420	-7.98045205741543\\
-410	-7.86973197872509\\
-400	-7.75900853528909\\
-390	-7.64828172700496\\
-380	-7.53755155377063\\
-370	-7.42681801548363\\
-360	-7.31608111204191\\
-350	-7.20534084334341\\
-340	-7.09459720928565\\
-330	-6.98385020976576\\
-320	-6.87309984468249\\
-310	-6.76234611393296\\
-300	-6.65158901741513\\
-290	-6.54082855502651\\
-280	-6.43006472666464\\
-270	-6.31929753222786\\
-260	-6.20852697161331\\
-250	-6.09775304471891\\
-240	-5.98697575144261\\
-230	-5.87619509168154\\
-220	-5.76541106533362\\
-210	-5.65462367229681\\
-200	-5.54383291246862\\
-190	-5.4330387857466\\
-180	-5.32224129202869\\
-170	-5.211440431212\\
-160	-5.10063620319488\\
-150	-4.98982860787487\\
-140	-4.8790176451495\\
-130	-4.76820331491629\\
-120	-4.6573856170732\\
-110	-4.54656455151774\\
-100	-4.43574011814786\\
-90	-4.32491231686069\\
-80	-4.21408114755376\\
-70	-4.1032466101254\\
-60	-3.99240870447316\\
-50	-3.88156743049417\\
-40	-3.77072278808635\\
-30	-3.65987477714764\\
-20	-3.54902339757518\\
-10	-3.4381686492665\\
0	-3.32731053211994\\
10	-3.21644904603262\\
20	-3.10558419090168\\
30	-2.99471596662586\\
40	-2.88384437310229\\
50	-2.7729694102281\\
60	-2.66209107790203\\
70	-2.55120937602041\\
80	-2.44032430448117\\
90	-2.32943586318224\\
100	-2.21854405202157\\
110	-2.10764887089547\\
120	-1.9967503197027\\
130	-1.88584839834037\\
140	-1.77494310670643\\
150	-1.664034444698\\
160	-1.55312241221303\\
170	-1.44220700914863\\
180	-1.33128823540275\\
190	-1.22036609087333\\
200	-1.10944057545667\\
210	-0.99851168905154\\
220	-0.887579431555054\\
230	-0.776643802864342\\
240	-0.665704802878153\\
250	-0.554762431492804\\
260	-0.443816688606233\\
270	-0.332867574116378\\
280	-0.221915087920365\\
290	-0.110959229916133\\
300	0\\
310	0.110962601928478\\
320	0.221928575973791\\
330	0.332897922237194\\
340	0.443870640821558\\
350	0.554846731828945\\
360	0.665826195362229\\
370	0.77680903152428\\
380	0.887795240417162\\
390	0.998784822143747\\
400	1.10977777680529\\
410	1.22077410450547\\
420	1.33177380534634\\
430	1.44277687943161\\
440	1.55378332686169\\
450	1.66479314774028\\
460	1.77580634217026\\
470	1.88682291025286\\
480	1.99784285209259\\
490	2.10886616778988\\
500	2.21989285744843\\
510	2.33092292117028\\
520	2.44195635905832\\
530	2.55299317121542\\
540	2.66403335774363\\
550	2.77507691874503\\
560	2.88612385432329\\
570	2.99717416457966\\
580	3.10822784961783\\
590	3.21928490953987\\
600	3.33034534444863\\
610	3.44140915444619\\
620	3.55247633963542\\
630	3.66354690011838\\
640	3.77462083599875\\
650	3.88569814737859\\
660	3.99677883435997\\
670	4.10786289704576\\
680	4.21895033553884\\
690	4.33004114994125\\
700	4.44113534035669\\
710	4.55223290688642\\
720	4.66333384963329\\
730	4.774438168701\\
740	4.88554586419079\\
750	4.99665693620635\\
760	5.10777138484974\\
770	5.21888921022303\\
780	5.33001041242989\\
790	5.44113499157319\\
800	5.5522629477542\\
810	5.66339428107578\\
820	5.77452899164162\\
830	5.88566707955377\\
840	5.9968085449143\\
850	6.1079533878269\\
860	6.2191016083928\\
870	6.33025320671652\\
880	6.44140818289929\\
890	6.552566537044\\
900	6.66372826925351\\
910	6.7748933796307\\
920	6.88606186827843\\
930	6.99723373529878\\
940	7.1084089807938\\
950	7.21958760486717\\
960	7.33076960762178\\
970	7.44195498915967\\
980	7.55314374958373\\
990	7.66433588899601\\
1000	7.77553140750101\\
1010	7.88673030519917\\
1020	7.99793258219417\\
1030	8.10913823858969\\
1040	8.220347274487\\
1050	8.33155968998895\\
1060	8.44277548519923\\
1070	8.5539946602191\\
1080	8.66521721515224\\
1090	8.77644315010151\\
1100	8.8876724651698\\
1110	8.99890516045835\\
1120	9.11014123607166\\
1130	9.22138069211178\\
1140	9.33262352868078\\
1150	9.44386974588235\\
1160	9.55511934381855\\
1170	9.66637232259224\\
1180	9.7776286823063\\
1190	9.88888842306442\\
1200	10.0001515449678\\
1210	10.1114180481194\\
1220	10.2226879326229\\
1230	10.3339611985803\\
1240	10.4452378460944\\
1250	10.5565178752691\\
1260	10.6678012862055\\
1270	10.7790880790073\\
1280	10.8903782537773\\
1290	11.0016718106178\\
1300	11.1129687496322\\
1310	11.2242690709219\\
1320	11.3355727745923\\
1330	11.4468798607436\\
1340	11.5581903294797\\
1350	11.6695041809033\\
1360	11.7808214151174\\
1370	11.892142032224\\
1380	12.0034660323268\\
1390	12.1147934155286\\
1400	12.2261241819316\\
1410	12.3374583316395\\
1420	12.4487958647541\\
1430	12.5601367813786\\
1440	12.6714810816165\\
1450	12.782828765569\\
1460	12.8941798333408\\
1470	13.0055342850337\\
1480	13.1168921207508\\
1490	13.2282533405941\\
1500	13.339617944668\\
1510	13.4509859330747\\
1520	13.5623573059161\\
1530	13.673732063296\\
1540	13.7851102053181\\
1550	13.8964917320836\\
1560	14.0078766436954\\
1570	14.1192649402579\\
1580	14.2306566218725\\
1590	14.3420516886427\\
1600	14.4534501406716\\
1610	14.5648519780619\\
1620	14.6762572009157\\
1630	14.7876658093375\\
1640	14.8990778034286\\
1650	15.0104931832926\\
1660	15.1219119490317\\
1670	15.2333341007502\\
1680	15.3447596385503\\
1690	15.456188562534\\
1700	15.5676208728058\\
1710	15.679056569467\\
1720	15.7904956526221\\
1730	15.901938122373\\
1740	16.0133839788228\\
1750	16.124833222075\\
1760	16.236285852231\\
1770	16.3477418693952\\
1780	16.4592012736705\\
1790	16.5706640651589\\
1800	16.6821302439642\\
1810	16.7935998101883\\
1820	16.9050727639351\\
1830	17.0165491053073\\
1840	17.1280288344087\\
1850	17.2395119513404\\
1860	17.3509984562062\\
1870	17.4624883491098\\
1880	17.5739816301531\\
1890	17.6854782994392\\
1900	17.7969783570717\\
1910	17.9084818031533\\
1920	18.0199886377871\\
1930	18.1314988610759\\
1940	18.2430124731226\\
1950	18.35452947403\\
1960	18.4660498639009\\
1970	18.5775736428392\\
1980	18.6891008109477\\
1990	18.8006313683291\\
2000	18.9121653150865\\
2010	19.0237026513227\\
2020	19.1352433771405\\
2030	19.2467874926428\\
2040	19.3583349979341\\
2050	19.4698858931157\\
2060	19.5814401782913\\
2070	19.6929978535645\\
2080	19.8045589190374\\
2090	19.9161233748128\\
2100	20.0276912209945\\
2110	20.1392624576862\\
2120	20.250837084989\\
2130	20.3624151030075\\
2140	20.4739965118437\\
2150	20.5855813116013\\
2160	20.6971695023841\\
2170	20.8087610842932\\
2180	20.9203560574332\\
2190	21.0319544219068\\
2200	21.1435561778171\\
2210	21.255161325266\\
2220	21.3667698643589\\
2230	21.478381795197\\
2240	21.5899971178832\\
2250	21.701615832522\\
2260	21.8132379392162\\
2270	21.9248634380679\\
2280	22.0364923291808\\
2290	22.1481246126577\\
2300	22.2597602886024\\
2310	22.3713993571178\\
2320	22.4830418183058\\
2330	22.594687672271\\
2340	22.7063369191162\\
2350	22.8179895589435\\
2360	22.9296455918575\\
2370	23.0413050179601\\
2380	23.152967837355\\
2390	23.2646340501452\\
2400	23.3763036564343\\
2410	23.4879766563243\\
2420	23.5996530499198\\
2430	23.711332837322\\
2440	23.8230160186362\\
2450	23.9347025939637\\
2460	24.0463925634089\\
2470	24.1580859270747\\
2480	24.2697826850633\\
2490	24.381482837479\\
2500	24.4931863844248\\
2510	24.6048933260035\\
2520	24.716603662318\\
2530	24.828317393472\\
2540	24.9400345195691\\
2550	25.0517550407115\\
2560	25.163478957002\\
2570	25.2752062685458\\
2580	25.3869369754435\\
2590	25.4986710778004\\
2600	25.6104085757193\\
2610	25.7221494693023\\
2620	25.833893758653\\
2630	25.9456414438751\\
2640	26.0573925250716\\
2650	26.1691470023461\\
2660	26.2809048758007\\
2670	26.3926661455398\\
2680	26.5044308116655\\
2690	26.6161988742823\\
2700	26.7279703334923\\
2710	26.8397451893992\\
2720	26.9515234421058\\
2730	27.0633050917166\\
2740	27.1750901383329\\
2750	27.2868785820599\\
2760	27.398670422999\\
2770	27.5104656612554\\
2780	27.6222642969304\\
2790	27.7340663301284\\
2800	27.8458717609524\\
2810	27.9576805895053\\
2820	28.0694928158906\\
2830	28.1813084402121\\
2840	28.2931274625727\\
2850	28.4049498830752\\
2860	28.5167757018233\\
2870	28.6286049189207\\
2880	28.7404375344694\\
2890	28.852273548574\\
2900	28.9641129613373\\
2910	29.0759557728622\\
2920	29.1878019832524\\
2930	29.2996515926107\\
2940	29.4115046010417\\
2950	29.5233610086474\\
2960	29.6352208155314\\
2970	29.7470840217968\\
2980	29.8589506275478\\
2990	29.9708206328867\\
3000	30.0826940379171\\
};
\addlegendentry{Exponential}

\addplot[only marks, mark=o, mark options={}, mark size=1.5000pt, draw=black] table[row sep=crcr]{%
x	y\\
239.61	3\\
307.7	3.5\\
330.14	3.67\\
374.64	4\\
427.39	4.4\\
590.07	5.67\\
737.43	6.86\\
800	7.38\\
913.51	8.33\\
1002.31	9.1\\
1094.21	9.91\\
1178.66	10.67\\
1274.04	11.54\\
1399.82	12.71\\
1478.67	13.47\\
1546.68	14.13\\
1557.98	14.24\\
1641.31	15.06\\
1688.54	15.53\\
1740.59	16.05\\
1805.89	16.71\\
1860.19	17.27\\
1938.41	18.09\\
1961.86	18.33\\
2017.31	18.92\\
2057.25	19.35\\
2114.56	19.96\\
2164	20.5\\
2200.24	20.9\\
2260.97	21.57\\
};
\addlegendentry{Calculated data}

\end{axis}
\end{tikzpicture}%
\end{figure}

In the graph,all data points of $T$ calculated from resistance are ploted against $R$.
The solid line represents the best fit quadratic curve which is the same as the one listed above.
The data points match up because this is the relation used to calculate the values of temperature.
The dotted line shows the best fit exponential curve.

From the figure we see that the quadratic curve almost indistinguishable from the exponential near the higher values,
however at lower values of $y$ both curves begin to diverge.
For the same value of $y$ the value of $x$ for the quadratic function is lower than the value of $x$ for the exponential.
This means that for lower values of resistance, the quadratic formula for temperature gives a lower than expected value
which is most likely the reason that the first value of calculated temperature is lower than the ambient temperature

For this reason, I shall be excluding all values lower than the given value of $800K$ of glowing temperature.

The second question in the points to ponder asks how much error is in neglecting the $T_a$ term in the heat exchange expression.
It is safe to assume that the ambient temperature ranges between $300K$ to $400K$.(input source)

For this we see the following table
\begin{table}[!ht]
    \centering
    \begin{tabular}{|c|c|}
    \hline
        $T_a$ & $T^4-T_a^4$ \\ \hline
        0 & 4.10E+11 \\ \hline
        20 & 4.10E+11 \\ \hline
        40 & 4.10E+11 \\ \hline
        60 & 4.10E+11 \\ \hline
        80 & 4.10E+11 \\ \hline
        100 & 4.10E+11 \\ \hline
        120 & 4.09E+11 \\ \hline
        140 & 4.09E+11 \\ \hline
        160 & 4.09E+11 \\ \hline
        180 & 4.09E+11 \\ \hline
        200 & 4.08E+11 \\ \hline
        220 & 4.07E+11 \\ \hline
        240 & 4.06E+11 \\ \hline
        260 & 4.05E+11 \\ \hline
        280 & 4.03E+11 \\ \hline
        300 & 4.02E+11 \\ \hline
        320 & 3.99E+11 \\ \hline
        340 & 3.96E+11 \\ \hline
        360 & 3.93E+11 \\ \hline
        380 & 3.89E+11 \\ \hline
        400 & 3.84E+11 \\ \hline
    \end{tabular}
    \caption{Values are calculated for $T$=800$K$}
\end{table}

\noindent From the table we can calculate the worst case error in  neglecting the ambient temperature
\[\frac{4.10-3.84}{3.84}=6.77\%\]
As the temperature increases, due to the magnitude of $T^4$ the error becomes small.

%Realistically, the actual errors for neglecting ambient temperature is as given in the below figure(credit)

I shall not be ignoring the ambient temperature completely, but considering it to be constant at arount $300K$.
\begin{figure}[!ht]
    \centering
    \setlength\fwidth{0.75\textwidth}
    \setlength\fheight{0.45\textwidth}
    % This file was created by matlab2tikz.
%
\definecolor{mycolor1}{rgb}{0.85000,0.32500,0.09800}%
\definecolor{mycolor2}{rgb}{0.12941,0.12941,0.12941}%
%
\begin{tikzpicture}

\begin{axis}[%
width=0.951\fwidth,
height=\fheight,
at={(0\fwidth,0\fheight)},
scale only axis,
xmin=700,
xmax=2400,
xlabel style={font=\color{mycolor2}},
xlabel={Temperature},
ymin=0.01,
ymax=2.5,
ylabel style={font=\color{mycolor2}},
ylabel={Power},
axis background/.style={fill=white},
axis x line*=bottom,
axis y line*=left,
xmajorgrids,
ymajorgrids,
legend style={at={(0.03,0.97)}, anchor=north west, legend cell align=left, align=left}
]
\addplot[only marks, mark=o, mark options={}, mark size=1.5000pt, draw=black] table[row sep=crcr]{%
x	y\\
800	0.0472\\
913.51	0.0675\\
1002.31	0.091\\
1094.21	0.1199\\
1178.66	0.1536\\
1274.04	0.195\\
1399.82	0.2492\\
1478.67	0.303\\
1546.68	0.3616\\
1557.98	0.4114\\
1641.31	0.4878\\
1688.54	0.5605\\
1740.59	0.642\\
1805.89	0.7371\\
1860.19	0.836\\
1938.41	0.9568\\
1961.86	1.056\\
2017.31	1.1825\\
2057.25	1.3078\\
2114.56	1.4553\\
2164	1.6072\\
2200.24	1.7574\\
2260.97	1.941\\
};
\addlegendentry{Data}

\addplot [color=mycolor1]
  table[row sep=crcr]{%
600	0.00635604282181177\\
605	0.0065979883693965\\
610	0.00684656756842957\\
615	0.00710190568316747\\
620	0.00736412929169905\\
625	0.00763336628887605\\
630	0.00790974588922637\\
635	0.00819339862985033\\
640	0.0084844563733001\\
645	0.00878305231044257\\
650	0.00908932096330589\\
655	0.00940339818790987\\
660	0.00972542117708046\\
665	0.0100555284632486\\
670	0.0103938599212334\\
675	0.0107405567710105\\
680	0.0110957615804644\\
685	0.011459618268127\\
690	0.0118322721059009\\
695	0.0122138697217679\\
700	0.0126045591024837\\
705	0.0130044895962579\\
710	0.013413811915421\\
715	0.013832678139076\\
720	0.0142612417157381\\
725	0.0146996574659594\\
730	0.0151480815849415\\
735	0.0156066716451341\\
740	0.016075586598821\\
745	0.0165549867806933\\
750	0.0170450339104095\\
755	0.0175458910951437\\
760	0.0180577228321206\\
765	0.018580695011139\\
770	0.0191149749170827\\
775	0.0196607312324196\\
780	0.0202181340396888\\
785	0.0207873548239764\\
790	0.021368566475379\\
795	0.0219619432914566\\
800	0.0225676609796731\\
805	0.023185896659827\\
810	0.0238168288664696\\
815	0.0244606375513133\\
820	0.0251175040856281\\
825	0.0257876112626284\\
830	0.0264711432998485\\
835	0.0271682858415076\\
840	0.0278792259608645\\
845	0.0286041521625625\\
850	0.0293432543849632\\
855	0.0300967240024708\\
860	0.0308647538278465\\
865	0.0316475381145127\\
870	0.0324452725588475\\
875	0.0332581543024698\\
880	0.0340863819345143\\
885	0.0349301554938976\\
890	0.0357896764715744\\
895	0.0366651478127847\\
900	0.0375567739192917\\
905	0.0384647606516106\\
910	0.0393893153312283\\
915	0.0403306467428143\\
920	0.0412889651364224\\
925	0.0422644822296841\\
930	0.0432574112099929\\
935	0.0442679667366803\\
940	0.0452963649431828\\
945	0.0463428234392013\\
950	0.0474075613128509\\
955	0.0484907991328035\\
960	0.0495927589504219\\
965	0.0507136643018849\\
970	0.0518537402103059\\
975	0.0530132131878419\\
980	0.0541923112377957\\
985	0.0553912638567097\\
990	0.056610302036452\\
995	0.0578496582662949\\
1000	0.0591095665349853\\
1005	0.0603902623328084\\
1010	0.061691982653643\\
1015	0.0630149659970097\\
1020	0.0643594523701117\\
1025	0.0657256832898683\\
1030	0.067113901784941\\
1035	0.0685243523977523\\
1040	0.0699572811864982\\
1045	0.0714129357271519\\
1050	0.0728915651154624\\
1055	0.0743934199689446\\
1060	0.0759187524288637\\
1065	0.0774678161622116\\
1070	0.0790408663636774\\
1075	0.0806381597576106\\
1080	0.0822599545999781\\
1085	0.083906510680314\\
1090	0.0855780893236634\\
1095	0.087274953392519\\
1100	0.0889973672887517\\
1105	0.0907455969555351\\
1110	0.0925199098792626\\
1115	0.094320575091459\\
1120	0.0961478631706856\\
1125	0.0980020462444392\\
1130	0.0998833979910441\\
1135	0.10179219364154\\
1140	0.10372870998156\\
1145	0.105693225353209\\
1150	0.107686019656928\\
1155	0.109707374353358\\
1160	0.111757572465198\\
1165	0.113836898579054\\
1170	0.115945638847285\\
1175	0.11808408098984\\
1180	0.120252514296095\\
1185	0.122451229626676\\
1190	0.124680519415288\\
1195	0.126940677670522\\
1200	0.129231999977674\\
1205	0.131554783500548\\
1210	0.133909326983253\\
1215	0.136295930752\\
1220	0.138714896716894\\
1225	0.141166528373709\\
1230	0.143651130805677\\
1235	0.146169010685252\\
1240	0.148720476275882\\
1245	0.151305837433774\\
1250	0.153925405609644\\
1255	0.156579493850477\\
1260	0.159268416801271\\
1265	0.161992490706776\\
1270	0.164752033413237\\
1275	0.16754736437012\\
1280	0.170378804631845\\
1285	0.173246676859504\\
1290	0.176151305322578\\
1295	0.179093015900652\\
1300	0.18207213608512\\
1305	0.18508899498089\\
1310	0.188143923308081\\
1315	0.191237253403714\\
1320	0.194369319223405\\
1325	0.197540456343043\\
1330	0.200751001960474\\
1335	0.204001294897176\\
1340	0.207291675599922\\
1345	0.210622486142455\\
1350	0.21399407022714\\
1355	0.217406773186626\\
1360	0.220860941985498\\
1365	0.224356925221919\\
1370	0.227895073129279\\
1375	0.231475737577831\\
1380	0.235099272076325\\
1385	0.238766031773638\\
1390	0.242476373460403\\
1395	0.246230655570625\\
1400	0.250029238183299\\
1405	0.253872483024029\\
1410	0.257760753466628\\
1415	0.26169441453473\\
1420	0.265673832903384\\
1425	0.269699376900656\\
1430	0.273771416509217\\
1435	0.277890323367933\\
1440	0.282056470773448\\
1445	0.286270233681767\\
1450	0.290531988709826\\
1455	0.29484211413707\\
1460	0.299200989907019\\
1465	0.303608997628831\\
1470	0.308066520578864\\
1475	0.31257394370223\\
1480	0.317131653614351\\
1485	0.321740038602505\\
1490	0.326399488627372\\
1495	0.331110395324575\\
1500	0.335873152006217\\
1505	0.340688153662413\\
1510	0.345555796962825\\
1515	0.350476480258181\\
1520	0.355450603581805\\
1525	0.360478568651128\\
1530	0.36556077886921\\
1535	0.370697639326248\\
1540	0.375889556801084\\
1545	0.381136939762711\\
1550	0.386440198371772\\
1555	0.391799744482058\\
1560	0.397215991642002\\
1565	0.40268935509617\\
1570	0.408220251786746\\
1575	0.413809100355016\\
1580	0.419456321142849\\
1585	0.425162336194171\\
1590	0.430927569256442\\
1595	0.43675244578212\\
1600	0.442637392930132\\
1605	0.448582839567338\\
1610	0.454589216269982\\
1615	0.460656955325161\\
1620	0.466786490732267\\
1625	0.472978258204443\\
1630	0.479232695170026\\
1635	0.485550240773994\\
1640	0.491931335879401\\
1645	0.498376423068819\\
1650	0.504885946645767\\
1655	0.511460352636147\\
1660	0.518100088789665\\
1665	0.524805604581257\\
1670	0.531577351212514\\
1675	0.538415781613096\\
1680	0.545321350442146\\
1685	0.552294514089706\\
1690	0.55933573067812\\
1695	0.566445460063446\\
1700	0.573624163836854\\
1705	0.580872305326025\\
1710	0.588190349596552\\
1715	0.59557876345333\\
1720	0.603038015441948\\
1725	0.610568575850076\\
1730	0.618170916708849\\
1735	0.625845511794254\\
1740	0.633592836628499\\
1745	0.641413368481399\\
1750	0.649307586371743\\
1755	0.657275971068667\\
1760	0.665319005093021\\
1765	0.673437172718732\\
1770	0.681630959974168\\
1775	0.689900854643495\\
1780	0.698247346268034\\
1785	0.706670926147615\\
1790	0.715172087341927\\
1795	0.723751324671865\\
1800	0.732409134720877\\
1805	0.741146015836304\\
1810	0.749962468130721\\
1815	0.758858993483276\\
1820	0.767836095541019\\
1825	0.776894279720239\\
1830	0.786034053207789\\
1835	0.795255924962415\\
1840	0.804560405716078\\
1845	0.813948007975275\\
1850	0.823419246022359\\
1855	0.832974635916853\\
1860	0.842614695496764\\
1865	0.852339944379892\\
1870	0.862150903965143\\
1875	0.872048097433826\\
1880	0.882032049750966\\
1885	0.892103287666595\\
1890	0.902262339717058\\
1895	0.912509736226302\\
1900	0.922846009307172\\
1905	0.933271692862702\\
1910	0.943787322587401\\
1915	0.95439343596854\\
1920	0.965090572287433\\
1925	0.975879272620721\\
1930	0.986760079841646\\
1935	0.99773353862133\\
1940	1.00880019543005\\
1945	1.0199605985385\\
1950	1.03121529801906\\
1955	1.0425648457471\\
1960	1.05400979540217\\
1965	1.06555070246933\\
1970	1.07718812424036\\
1975	1.08892261981507\\
1980	1.1007547501025\\
1985	1.1126850778222\\
1990	1.12471416750548\\
1995	1.13684258549665\\
2000	1.14907089995427\\
2005	1.16139968085239\\
2010	1.17382949998179\\
2015	1.18636093095122\\
2020	1.19899454918865\\
2025	1.21173093194247\\
2030	1.22457065828275\\
2035	1.23751430910248\\
2040	1.25056246711875\\
2045	1.26371571687402\\
2050	1.27697464473733\\
2055	1.29033983890552\\
2060	1.30381188940443\\
2065	1.31739138809015\\
2070	1.33107892865022\\
2075	1.34487510660483\\
2080	1.35878051930803\\
2085	1.37279576594898\\
2090	1.38692144755309\\
2095	1.40115816698327\\
2100	1.41550652894114\\
2105	1.42996713996817\\
2110	1.44454060844695\\
2115	1.45922754460232\\
2120	1.47402856050263\\
2125	1.48894427006087\\
2130	1.50397528903589\\
2135	1.51912223503359\\
2140	1.53438572750808\\
2145	1.54976638776291\\
2150	1.56526483895219\\
2155	1.58088170608183\\
2160	1.59661761601066\\
2165	1.61247319745166\\
2170	1.62844908097307\\
2175	1.64454589899965\\
2180	1.66076428581373\\
2185	1.67710487755651\\
2190	1.69356831222912\\
2195	1.71015522969382\\
2200	1.72686627167518\\
2205	1.74370208176123\\
2210	1.76066330540459\\
2215	1.77775058992364\\
2220	1.79496458450372\\
2225	1.81230594019821\\
2230	1.82977530992971\\
2235	1.8473733484912\\
2240	1.86510071254719\\
2245	1.88295806063483\\
2250	1.90094605316507\\
2255	1.91906535242382\\
2260	1.93731662257307\\
2265	1.95570052965202\\
2270	1.97421774157824\\
2275	1.99286892814876\\
2280	2.01165476104126\\
2285	2.03057591381516\\
2290	2.04963306191274\\
2295	2.06882688266031\\
2300	2.08815805526928\\
2305	2.10762726083734\\
2310	2.12723518234952\\
2315	2.14698250467936\\
2320	2.16686991458999\\
2325	2.18689810073529\\
2330	2.20706775366094\\
2335	2.2273795658056\\
2340	2.24783423150196\\
2345	2.26843244697789\\
2350	2.28917491035755\\
2355	2.31006232166244\\
2360	2.33109538281257\\
2365	2.35227479762752\\
2370	2.37360127182756\\
2375	2.39507551303475\\
2380	2.416698230774\\
2385	2.43847013647421\\
2390	2.46039194346935\\
2395	2.48246436699954\\
2400	2.50468812421215\\
};
\addlegendentry{Best fit}

\end{axis}
\end{tikzpicture}%
\end{figure}

\noindent Performing regression for the formula
$y=c(x^a-300^a)$
we get the value of $a$ as 4.27 (with $R^2 = 0.9987$)

\begin{figure}[!ht]
    \centering
    \setlength\fwidth{0.75\textwidth}
    \setlength\fheight{0.45\textwidth}
    % This file was created by matlab2tikz.
%
\definecolor{mycolor1}{rgb}{0.85000,0.32500,0.09800}%
\definecolor{mycolor2}{rgb}{0.92900,0.69400,0.12500}%
\definecolor{mycolor3}{rgb}{0.12941,0.12941,0.12941}%
%
\begin{tikzpicture}

\begin{axis}[%
width=0.951\fwidth,
height=\fheight,
at={(0\fwidth,0\fheight)},
scale only axis,
xmode=log,
xmin=700,
xmax=2400,
xminorticks=true,
xlabel style={font=\color{mycolor3}},
xlabel={Temperature},
ymode=log,
ymin=0.01,
ymax=3,
yminorticks=true,
ylabel style={font=\color{mycolor3}},
ylabel={Power},
axis background/.style={fill=white},
legend style={at={(0.03,0.97)}, anchor=north west, legend cell align=left, align=left}
]
\addplot[only marks, mark=o, mark options={}, mark size=1.5000pt, draw=black] table[row sep=crcr]{%
x	y\\
800	0.0472\\
913.51	0.0675\\
1002.31	0.091\\
1094.21	0.1199\\
1178.66	0.1536\\
1274.04	0.195\\
1399.82	0.2492\\
1478.67	0.303\\
1546.68	0.3616\\
1557.98	0.4114\\
1641.31	0.4878\\
1688.54	0.5605\\
1740.59	0.642\\
1805.89	0.7371\\
1860.19	0.836\\
1938.41	0.9568\\
1961.86	1.056\\
2017.31	1.1825\\
2057.25	1.3078\\
2114.56	1.4553\\
2164	1.6072\\
2200.24	1.7574\\
2260.97	1.941\\
};
\addlegendentry{Data}

\addplot [color=mycolor1]
  table[row sep=crcr]{%
600	0.00635604282181177\\
605	0.0065979883693965\\
610	0.00684656756842957\\
615	0.00710190568316747\\
620	0.00736412929169905\\
625	0.00763336628887605\\
630	0.00790974588922637\\
635	0.00819339862985033\\
640	0.0084844563733001\\
645	0.00878305231044257\\
650	0.00908932096330589\\
655	0.00940339818790987\\
660	0.00972542117708046\\
665	0.0100555284632486\\
670	0.0103938599212334\\
675	0.0107405567710105\\
680	0.0110957615804644\\
685	0.011459618268127\\
690	0.0118322721059009\\
695	0.0122138697217679\\
700	0.0126045591024837\\
705	0.0130044895962579\\
710	0.013413811915421\\
715	0.013832678139076\\
720	0.0142612417157381\\
725	0.0146996574659594\\
730	0.0151480815849415\\
735	0.0156066716451341\\
740	0.016075586598821\\
745	0.0165549867806933\\
750	0.0170450339104095\\
755	0.0175458910951437\\
760	0.0180577228321206\\
765	0.018580695011139\\
770	0.0191149749170827\\
775	0.0196607312324196\\
780	0.0202181340396888\\
785	0.0207873548239764\\
790	0.021368566475379\\
795	0.0219619432914566\\
800	0.0225676609796731\\
805	0.023185896659827\\
810	0.0238168288664696\\
815	0.0244606375513133\\
820	0.0251175040856281\\
825	0.0257876112626284\\
830	0.0264711432998485\\
835	0.0271682858415076\\
840	0.0278792259608645\\
845	0.0286041521625625\\
850	0.0293432543849632\\
855	0.0300967240024708\\
860	0.0308647538278465\\
865	0.0316475381145127\\
870	0.0324452725588475\\
875	0.0332581543024698\\
880	0.0340863819345143\\
885	0.0349301554938976\\
890	0.0357896764715744\\
895	0.0366651478127847\\
900	0.0375567739192917\\
905	0.0384647606516106\\
910	0.0393893153312283\\
915	0.0403306467428143\\
920	0.0412889651364224\\
925	0.0422644822296841\\
930	0.0432574112099929\\
935	0.0442679667366803\\
940	0.0452963649431828\\
945	0.0463428234392013\\
950	0.0474075613128509\\
955	0.0484907991328035\\
960	0.0495927589504219\\
965	0.0507136643018849\\
970	0.0518537402103059\\
975	0.0530132131878419\\
980	0.0541923112377957\\
985	0.0553912638567097\\
990	0.056610302036452\\
995	0.0578496582662949\\
1000	0.0591095665349853\\
1005	0.0603902623328084\\
1010	0.061691982653643\\
1015	0.0630149659970097\\
1020	0.0643594523701117\\
1025	0.0657256832898683\\
1030	0.067113901784941\\
1035	0.0685243523977523\\
1040	0.0699572811864982\\
1045	0.0714129357271519\\
1050	0.0728915651154624\\
1055	0.0743934199689446\\
1060	0.0759187524288637\\
1065	0.0774678161622116\\
1070	0.0790408663636774\\
1075	0.0806381597576106\\
1080	0.0822599545999781\\
1085	0.083906510680314\\
1090	0.0855780893236634\\
1095	0.087274953392519\\
1100	0.0889973672887517\\
1105	0.0907455969555351\\
1110	0.0925199098792626\\
1115	0.094320575091459\\
1120	0.0961478631706856\\
1125	0.0980020462444392\\
1130	0.0998833979910441\\
1135	0.10179219364154\\
1140	0.10372870998156\\
1145	0.105693225353209\\
1150	0.107686019656928\\
1155	0.109707374353358\\
1160	0.111757572465198\\
1165	0.113836898579054\\
1170	0.115945638847285\\
1175	0.11808408098984\\
1180	0.120252514296095\\
1185	0.122451229626676\\
1190	0.124680519415288\\
1195	0.126940677670522\\
1200	0.129231999977674\\
1205	0.131554783500548\\
1210	0.133909326983253\\
1215	0.136295930752\\
1220	0.138714896716894\\
1225	0.141166528373709\\
1230	0.143651130805677\\
1235	0.146169010685252\\
1240	0.148720476275882\\
1245	0.151305837433774\\
1250	0.153925405609644\\
1255	0.156579493850477\\
1260	0.159268416801271\\
1265	0.161992490706776\\
1270	0.164752033413237\\
1275	0.16754736437012\\
1280	0.170378804631845\\
1285	0.173246676859504\\
1290	0.176151305322578\\
1295	0.179093015900652\\
1300	0.18207213608512\\
1305	0.18508899498089\\
1310	0.188143923308081\\
1315	0.191237253403714\\
1320	0.194369319223405\\
1325	0.197540456343043\\
1330	0.200751001960474\\
1335	0.204001294897176\\
1340	0.207291675599922\\
1345	0.210622486142455\\
1350	0.21399407022714\\
1355	0.217406773186626\\
1360	0.220860941985498\\
1365	0.224356925221919\\
1370	0.227895073129279\\
1375	0.231475737577831\\
1380	0.235099272076325\\
1385	0.238766031773638\\
1390	0.242476373460403\\
1395	0.246230655570625\\
1400	0.250029238183299\\
1405	0.253872483024029\\
1410	0.257760753466628\\
1415	0.26169441453473\\
1420	0.265673832903384\\
1425	0.269699376900656\\
1430	0.273771416509217\\
1435	0.277890323367933\\
1440	0.282056470773448\\
1445	0.286270233681767\\
1450	0.290531988709826\\
1455	0.29484211413707\\
1460	0.299200989907019\\
1465	0.303608997628831\\
1470	0.308066520578864\\
1475	0.31257394370223\\
1480	0.317131653614351\\
1485	0.321740038602505\\
1490	0.326399488627372\\
1495	0.331110395324575\\
1500	0.335873152006217\\
1505	0.340688153662413\\
1510	0.345555796962825\\
1515	0.350476480258181\\
1520	0.355450603581805\\
1525	0.360478568651128\\
1530	0.36556077886921\\
1535	0.370697639326248\\
1540	0.375889556801084\\
1545	0.381136939762711\\
1550	0.386440198371772\\
1555	0.391799744482058\\
1560	0.397215991642002\\
1565	0.40268935509617\\
1570	0.408220251786746\\
1575	0.413809100355016\\
1580	0.419456321142849\\
1585	0.425162336194171\\
1590	0.430927569256442\\
1595	0.43675244578212\\
1600	0.442637392930132\\
1605	0.448582839567338\\
1610	0.454589216269982\\
1615	0.460656955325161\\
1620	0.466786490732267\\
1625	0.472978258204443\\
1630	0.479232695170026\\
1635	0.485550240773994\\
1640	0.491931335879401\\
1645	0.498376423068819\\
1650	0.504885946645767\\
1655	0.511460352636147\\
1660	0.518100088789665\\
1665	0.524805604581257\\
1670	0.531577351212514\\
1675	0.538415781613096\\
1680	0.545321350442146\\
1685	0.552294514089706\\
1690	0.55933573067812\\
1695	0.566445460063446\\
1700	0.573624163836854\\
1705	0.580872305326025\\
1710	0.588190349596552\\
1715	0.59557876345333\\
1720	0.603038015441948\\
1725	0.610568575850076\\
1730	0.618170916708849\\
1735	0.625845511794254\\
1740	0.633592836628499\\
1745	0.641413368481399\\
1750	0.649307586371743\\
1755	0.657275971068667\\
1760	0.665319005093021\\
1765	0.673437172718732\\
1770	0.681630959974168\\
1775	0.689900854643495\\
1780	0.698247346268034\\
1785	0.706670926147615\\
1790	0.715172087341927\\
1795	0.723751324671865\\
1800	0.732409134720877\\
1805	0.741146015836304\\
1810	0.749962468130721\\
1815	0.758858993483276\\
1820	0.767836095541019\\
1825	0.776894279720239\\
1830	0.786034053207789\\
1835	0.795255924962415\\
1840	0.804560405716078\\
1845	0.813948007975275\\
1850	0.823419246022359\\
1855	0.832974635916853\\
1860	0.842614695496764\\
1865	0.852339944379892\\
1870	0.862150903965143\\
1875	0.872048097433826\\
1880	0.882032049750966\\
1885	0.892103287666595\\
1890	0.902262339717058\\
1895	0.912509736226302\\
1900	0.922846009307172\\
1905	0.933271692862702\\
1910	0.943787322587401\\
1915	0.95439343596854\\
1920	0.965090572287433\\
1925	0.975879272620721\\
1930	0.986760079841646\\
1935	0.99773353862133\\
1940	1.00880019543005\\
1945	1.0199605985385\\
1950	1.03121529801906\\
1955	1.0425648457471\\
1960	1.05400979540217\\
1965	1.06555070246933\\
1970	1.07718812424036\\
1975	1.08892261981507\\
1980	1.1007547501025\\
1985	1.1126850778222\\
1990	1.12471416750548\\
1995	1.13684258549665\\
2000	1.14907089995427\\
2005	1.16139968085239\\
2010	1.17382949998179\\
2015	1.18636093095122\\
2020	1.19899454918865\\
2025	1.21173093194247\\
2030	1.22457065828275\\
2035	1.23751430910248\\
2040	1.25056246711875\\
2045	1.26371571687402\\
2050	1.27697464473733\\
2055	1.29033983890552\\
2060	1.30381188940443\\
2065	1.31739138809015\\
2070	1.33107892865022\\
2075	1.34487510660483\\
2080	1.35878051930803\\
2085	1.37279576594898\\
2090	1.38692144755309\\
2095	1.40115816698327\\
2100	1.41550652894114\\
2105	1.42996713996817\\
2110	1.44454060844695\\
2115	1.45922754460232\\
2120	1.47402856050263\\
2125	1.48894427006087\\
2130	1.50397528903589\\
2135	1.51912223503359\\
2140	1.53438572750808\\
2145	1.54976638776291\\
2150	1.56526483895219\\
2155	1.58088170608183\\
2160	1.59661761601066\\
2165	1.61247319745166\\
2170	1.62844908097307\\
2175	1.64454589899965\\
2180	1.66076428581373\\
2185	1.67710487755651\\
2190	1.69356831222912\\
2195	1.71015522969382\\
2200	1.72686627167518\\
2205	1.74370208176123\\
2210	1.76066330540459\\
2215	1.77775058992364\\
2220	1.79496458450372\\
2225	1.81230594019821\\
2230	1.82977530992971\\
2235	1.8473733484912\\
2240	1.86510071254719\\
2245	1.88295806063483\\
2250	1.90094605316507\\
2255	1.91906535242382\\
2260	1.93731662257307\\
2265	1.95570052965202\\
2270	1.97421774157824\\
2275	1.99286892814876\\
2280	2.01165476104126\\
2285	2.03057591381516\\
2290	2.04963306191274\\
2295	2.06882688266031\\
2300	2.08815805526928\\
2305	2.10762726083734\\
2310	2.12723518234952\\
2315	2.14698250467936\\
2320	2.16686991458999\\
2325	2.18689810073529\\
2330	2.20706775366094\\
2335	2.2273795658056\\
2340	2.24783423150196\\
2345	2.26843244697789\\
2350	2.28917491035755\\
2355	2.31006232166244\\
2360	2.33109538281257\\
2365	2.35227479762752\\
2370	2.37360127182756\\
2375	2.39507551303475\\
2380	2.416698230774\\
2385	2.43847013647421\\
2390	2.46039194346935\\
2395	2.48246436699954\\
2400	2.50468812421215\\
};
\addlegendentry{Best fit}

\addplot [color=mycolor2, dashed]
  table[row sep=crcr]{%
600	0.0088011684\\
605	0.0091180322389535\\
610	0.009442850360216\\
615	0.0097757547813135\\
620	0.010116878606336\\
625	0.0104663560259375\\
630	0.010824322317336\\
635	0.0111909138443135\\
640	0.011566268057216\\
645	0.0119505234929535\\
650	0.012343819775\\
655	0.0127462976133935\\
660	0.013158098804736\\
665	0.0135793662321935\\
670	0.014010243865496\\
675	0.0144508767609375\\
680	0.014901411061376\\
685	0.0153619939962335\\
690	0.015832773881496\\
695	0.0163139001197135\\
700	0.0168055232\\
705	0.0173077946980335\\
710	0.017820867276056\\
715	0.0183448946828735\\
720	0.018880031753856\\
725	0.0194264344109375\\
730	0.019984259662616\\
735	0.0205536656039535\\
740	0.021134811416576\\
745	0.0217278573686735\\
750	0.022332964815\\
755	0.0229502961968735\\
760	0.023580015042176\\
765	0.0242222859653535\\
770	0.024877274667416\\
775	0.0255451479359375\\
780	0.026226073645056\\
785	0.0269202207554735\\
790	0.027627759314456\\
795	0.0283488604558335\\
800	0.0290836964\\
805	0.0298324404539135\\
810	0.030595267011096\\
815	0.0313723515516335\\
820	0.032163870642176\\
825	0.0329700019359375\\
830	0.033790924172696\\
835	0.0346268171787935\\
840	0.035477861867136\\
845	0.0363442402371935\\
850	0.037226135375\\
855	0.0381237314531535\\
860	0.039037213730816\\
865	0.0399667685537135\\
870	0.040912583354136\\
875	0.0418748466509375\\
880	0.042853748049536\\
885	0.0438494782419135\\
890	0.044862229006616\\
895	0.0458921932087535\\
900	0.0469395648\\
905	0.0480045388185935\\
910	0.049087311389336\\
915	0.0501880797235935\\
920	0.051307042119296\\
925	0.0524443979609375\\
930	0.053600347719576\\
935	0.0547750929528335\\
940	0.055968836304896\\
945	0.0571817815065135\\
950	0.058414133375\\
955	0.0596660978142335\\
960	0.060937881814656\\
965	0.0622296934532735\\
970	0.063541741893656\\
975	0.0648742373859375\\
980	0.066227391266816\\
985	0.0676014159595535\\
990	0.068996524973976\\
995	0.0704129329064735\\
1000	0.07185085544\\
1005	0.0733105093440735\\
1010	0.074792112474776\\
1015	0.0762958837747535\\
1020	0.077822043273216\\
1025	0.0793708120859375\\
1030	0.080942412415256\\
1035	0.0825370675500735\\
1040	0.084155001865856\\
1045	0.0857964408246335\\
1050	0.087461610975\\
1055	0.0891507399521135\\
1060	0.090864056477696\\
1065	0.0926017903600335\\
1070	0.094364172493976\\
1075	0.0961514348609375\\
1080	0.097963810528896\\
1085	0.0998015336523935\\
1090	0.101664839472536\\
1095	0.103553964316993\\
1100	0.1054691456\\
1105	0.107410621822353\\
1110	0.109378632571416\\
1115	0.111373418521113\\
1120	0.113395221431936\\
1125	0.115444284150937\\
1130	0.117520850611736\\
1135	0.119625165834513\\
1140	0.121757475926016\\
1145	0.123918028079553\\
1150	0.126107070575\\
1155	0.128324852778793\\
1160	0.130571625143936\\
1165	0.132847639209993\\
1170	0.135153147603096\\
1175	0.137488404035937\\
1180	0.139853663307776\\
1185	0.142249181304433\\
1190	0.144675214998296\\
1195	0.147132022448313\\
1200	0.1496198628\\
1205	0.152138996285433\\
1210	0.154689684223256\\
1215	0.157272189018673\\
1220	0.159886774163456\\
1225	0.162533704235937\\
1230	0.165213244901016\\
1235	0.167925662910153\\
1240	0.170671226101376\\
1245	0.173450203399274\\
1250	0.176262864815\\
1255	0.179109481446273\\
1260	0.181990325477376\\
1265	0.184905670179153\\
1270	0.187855789909016\\
1275	0.190840960110937\\
1280	0.193861457315456\\
1285	0.196917559139673\\
1290	0.200009544287256\\
1295	0.203137692548433\\
1300	0.2063022848\\
1305	0.209503603005313\\
1310	0.212741930214296\\
1315	0.216017550563433\\
1320	0.219330749275776\\
1325	0.222681812660937\\
1330	0.226071028115096\\
1335	0.229498684120993\\
1340	0.232965070247936\\
1345	0.236470477151793\\
1350	0.240015196575\\
1355	0.243599521346553\\
1360	0.247223745382016\\
1365	0.250888163683513\\
1370	0.254593072339736\\
1375	0.258338768525937\\
1380	0.262125550503936\\
1385	0.265953717622114\\
1390	0.269823570315416\\
1395	0.273735410105353\\
1400	0.2776895396\\
1405	0.281686262493993\\
1410	0.285725883568536\\
1415	0.289808708691393\\
1420	0.293935044816896\\
1425	0.298105199985937\\
1430	0.302319483325976\\
1435	0.306578205051034\\
1440	0.310881676461696\\
1445	0.315230209945113\\
1450	0.319624118975\\
1455	0.324063718111633\\
1460	0.328549323001856\\
1465	0.333081250379073\\
1470	0.337659818063256\\
1475	0.342285344960937\\
1480	0.346958151065216\\
1485	0.351678557455753\\
1490	0.356446886298776\\
1495	0.361263460847074\\
1500	0.36612860544\\
1505	0.371042645503473\\
1510	0.376005907549976\\
1515	0.381018719178553\\
1520	0.386081409074816\\
1525	0.391194307010937\\
1530	0.396357743845656\\
1535	0.401572051524273\\
1540	0.406837563078656\\
1545	0.412154612627233\\
1550	0.417523535375\\
1555	0.422944667613513\\
1560	0.428418346720896\\
1565	0.433944911161833\\
1570	0.439524700487576\\
1575	0.445158055335937\\
1580	0.450845317431296\\
1585	0.456586829584593\\
1590	0.462382935693336\\
1595	0.468233980741593\\
1600	0.4741403108\\
1605	0.480102273025753\\
1610	0.486120215662616\\
1615	0.492194488040913\\
1620	0.498325440577536\\
1625	0.504513424775937\\
1630	0.510758793226136\\
1635	0.517061899604713\\
1640	0.523423098674816\\
1645	0.529842746286153\\
1650	0.536321199375\\
1655	0.542858815964193\\
1660	0.549455955163136\\
1665	0.556112977167794\\
1670	0.562830243260696\\
1675	0.569608115810938\\
1680	0.576446958274176\\
1685	0.583347135192633\\
1690	0.590309012195096\\
1695	0.597332955996913\\
1700	0.6044193344\\
1705	0.611568516292833\\
1710	0.618780871650456\\
1715	0.626056771534473\\
1720	0.633396588093056\\
1725	0.640800694560937\\
1730	0.648269465259416\\
1735	0.655803275596353\\
1740	0.663402502066176\\
1745	0.671067522249873\\
1750	0.678798714815\\
1755	0.686596459515673\\
1760	0.694461137192576\\
1765	0.702393129772953\\
1770	0.710392820270616\\
1775	0.718460592785937\\
1780	0.726596832505856\\
1785	0.734801925703873\\
1790	0.743076259740056\\
1795	0.751420223061033\\
1800	0.7598342052\\
1805	0.768318596776713\\
1810	0.776873789497496\\
1815	0.785500176155233\\
1820	0.794198150629376\\
1825	0.802968107885937\\
1830	0.811810443977496\\
1835	0.820725556043193\\
1840	0.829713842308736\\
1845	0.838775702086393\\
1850	0.847911535775\\
1855	0.857121744859953\\
1860	0.866406731913216\\
1865	0.875766900593313\\
1870	0.885202655645336\\
1875	0.894714402900937\\
1880	0.904302549278336\\
1885	0.913967502782313\\
1890	0.923709672504216\\
1895	0.933529468621953\\
1900	0.9434273024\\
1905	0.953403586189393\\
1910	0.963458733427736\\
1915	0.973593158639193\\
1920	0.983807277434496\\
1925	0.994101506510937\\
1930	1.00447626365238\\
1935	1.01493196772923\\
1940	1.0254690386985\\
1945	1.03608789760371\\
1950	1.046788966575\\
1955	1.05757266882903\\
1960	1.06843942866906\\
1965	1.07938967148487\\
1970	1.09042382375286\\
1975	1.10154231303594\\
1980	1.11274556798362\\
1985	1.12403401833195\\
1990	1.13540809490358\\
1995	1.14686822960767\\
2000	1.15841485544\\
2005	1.17004840648287\\
2010	1.18176931790518\\
2015	1.19357802596235\\
2020	1.20547496799642\\
2025	1.21746058243594\\
2030	1.22953530879606\\
2035	1.24169958767847\\
2040	1.25395386077146\\
2045	1.26629857084983\\
2050	1.278734161775\\
2055	1.29126107849491\\
2060	1.3038797670441\\
2065	1.31659067454363\\
2070	1.32939424920118\\
2075	1.34229094031094\\
2080	1.3552811982537\\
2085	1.36836547449679\\
2090	1.38154422159414\\
2095	1.39481789318619\\
2100	1.408186944\\
2105	1.42165182984915\\
2110	1.43521300763382\\
2115	1.44887093534071\\
2120	1.46262607204314\\
2125	1.47647887790094\\
2130	1.49042981416054\\
2135	1.50447934315491\\
2140	1.51862792830362\\
2145	1.53287603411275\\
2150	1.547224126175\\
2155	1.56167267116959\\
2160	1.57622213686234\\
2165	1.59087299210559\\
2170	1.6056257068383\\
2175	1.62048075208594\\
2180	1.63543859996058\\
2185	1.65049972366083\\
2190	1.6656645974719\\
2195	1.68093369676551\\
2200	1.696307498\\
2205	1.71178647872023\\
2210	1.72737111755766\\
2215	1.74306189423027\\
2220	1.75885928954266\\
2225	1.77476378538594\\
2230	1.79077586473782\\
2235	1.80689601166255\\
2240	1.82312471131098\\
2245	1.83946244992047\\
2250	1.855909714815\\
2255	1.87246699440507\\
2260	1.88913477818778\\
2265	1.90591355674675\\
2270	1.92280382175222\\
2275	1.93980606596094\\
2280	1.95692078321626\\
2285	1.97414846844807\\
2290	1.99148961767286\\
2295	2.00894472799363\\
2300	2.0265142976\\
2305	2.04419882576811\\
2310	2.0619988128607\\
2315	2.07991476032703\\
2320	2.09794717070298\\
2325	2.11609654761094\\
2330	2.1343633957599\\
2335	2.15274822094539\\
2340	2.17125153004954\\
2345	2.18987383104099\\
2350	2.208615632975\\
2355	2.22747744599335\\
2360	2.24645978132442\\
2365	2.26556315128311\\
2370	2.28478806927094\\
2375	2.30413504977594\\
2380	2.32360460837274\\
2385	2.34319726172251\\
2390	2.36291352757302\\
2395	2.38275392475855\\
2400	2.4027189732\\
};
\addlegendentry{power4}

\end{axis}
\end{tikzpicture}%
\end{figure}

When fitting is performed with the formula $y=c(x^4-300^4)$
we get the value of c to be $7.24376\times10^{-14}$

Now, $c = \epsilon \sigma A$.
The fillament can be thought of as a cylinder, thus the surface area is $2\pi r h$.
The radii($r$) and length($h$) of the filament were the order of $10^{-3}$(1mm) each.
$\epsilon$ lies between 0 and one.
This verifies that the order of magnitudes of calculated and given Stefan-Boltzmann constant are the same.
\[\frac{10^{-14}}{10^{-3}\times 10^{-3}} = 10^{-8}\]
There is not really much to comment about the value because we do not know the value of $\epsilon$ and $A$

\section*{Result}
We have verified that heat radiated is proportional to the fourth power of temperature.
\end{document}